\exercisesfor{4}

\begin{enumerate}
\def\labelenumi{\arabic{enumi}.}
\tightlist
\item
  Let G be a Lie group, H a normal Lie quasi-subgroup of G and \(\pi\) the canonical mapping of G onto G/H. There exists on G/H one and only one Lie group structure with the following property: for a homomorphism \(\theta\) of G/H into a Lie group \(G'\) to be a Lie group morphism, it is necessary and sufficient that \(\theta \circ \pi\) be a Lie group morphism. Moreover, L(G/H) is canonically isomorphic to L(G)/L(H). (Let Q be a Lie group germ such that L(Q) = L(G)/L(H). Show that, by shrinking Q if necessary, Q can be identified with an open neighbourhood of 0 in G/H and that Proposition 18 of §1 can then be applied to G/H.)
\end{enumerate}

¶ 2. Let G and H be two Lie groups and f a Lie group morphism of G into H.

\begin{enumerate}
\def\labelenumi{(\roman{enumi})}
\item
  The kernel \(\mathbf{N}\) of \(f\) is a normal Lie quasi-subgroup of \(\mathbf{G}\) and \(\mathbf{L}(\mathbf{N}) = \operatorname{Ker} \mathbf{L}(f)\) . (Use exponential mappings of \(\mathbf{G}\) and \(\mathbf{H}\) ).
\item
  Let \(g \colon \mathrm{G} / \mathrm{N} \to f(\mathrm{G})\) be the mapping derived from \(f\) when passing to the quotient. If \(f\) and \(\mathrm{L}(f)\) have closed images, and the topology of \(\mathrm{G}\) admits a countable base, then \(f(\mathrm{G})\) is a Lie quasi-subgroup of \(\mathrm{H}\) with Lie algebra \(\operatorname{Im} \mathrm{L}(f)\) and \(g\) is a Lie group isomorphism, where \(\mathrm{G} / \mathrm{N}\) has the structure defined in Exercise 1. (Reduce it by means of (i) to the case where \(\mathrm{N} = \{e\}\) .)
\end{enumerate}

\begin{enumerate}
\def\labelenumi{\arabic{enumi}.}
\setcounter{enumi}{2}
\tightlist
\item
  Let G be a Lie group, U an open neighbourhood of 0 in L(G) and \(\phi\) an analytic mapping of U into G such that \(\phi(0) = e\) and \(T_{0}\phi = \mathrm{Id}_{\mathrm{L}(G)}\) . The following conditions are equivalent:
\end{enumerate}

\begin{enumerate}
\def\labelenumi{(\alph{enumi})}
\item
  \(\phi\) is an exponential mapping;
\item
  There exists \(m \in \mathbf{Z}\) distinct from 0, 1, -1 such that \(\phi(mx) = \phi(x)^m\) in a neighbourhood of 0.
\end{enumerate}

(If condition (b) holds, show that \(\phi^{*}\) satisfies condition (iii) of Theorem 4.)

¶ 4. (a) Let K = R or C or an ultrametric field with residue field of characteristic \textgreater0. Let G be a Lie group over K, U an open neighbourhood of 0 in L(G) and \(\phi: U \to G\) a mapping which is differentiable at 0 and such
that \(\phi(0) = e\) , \(T_0(\phi) = \mathrm{Id}_{L(G)}\) and \(\phi((\lambda + \lambda')b) = \phi(\lambda b)\phi(\lambda'b)\) for \(\lambda b\) , \(\lambda'b\) , \((\lambda + \lambda')b\) in U. Then \(\phi\) coincides on a neighbourhood of 0 with an exponential mapping.

\begin{enumerate}
\def\labelenumi{(\alph{enumi})}
\setcounter{enumi}{1}
\tightlist
\item
  Let \(k\) be a field of characteristic 0. Let \(\mathbf{K}\) be the valued field \(k((\mathbf{X}))\) , which is of characteristic 0. Let \(\mathbf{G}\) be the additive Lie group \(k[[\mathbf{X}]]\) . Let \(\phi\) be the continuous \(k\) -linear mapping of \(\mathbf{G}\) into \(\mathbf{G}\) such that \(\phi(\mathbf{X}^n) = \mathbf{X}^n + \mathbf{X}^{2n}\) for every integer \(n \geqslant 0\) . Then \(\phi\) satisfies the conditions of (a), but coincides on no neighbourhood of 0 with an exponential mapping.
\end{enumerate}

\begin{enumerate}
\def\labelenumi{\arabic{enumi}.}
\setcounter{enumi}{4}
\tightlist
\item
  Let \((e_1, e_2, e_3)\) be the canonical basis of \(\mathbf{R}^3\) . We give \(\mathbf{R}^3\) the nilpotent Lie algebra structure \(\mathfrak{g}\) such that \([e_1, e_2] = e_3\) , \([e_1, e_3] = [e_2, e_3] = 0\) . Let \(c_{ijk}\) denote the constants of structure of \(\mathfrak{g}\) relative to the basis \((e_1, e_2, e_3)\) .
\end{enumerate}

\begin{enumerate}
\def\labelenumi{(\alph{enumi})}
\tightlist
\item
  Show that on \(\mathbf{R}^3\) the differential forms
\end{enumerate}

\[
\omega_ {1} = d x _ {1}, \quad \omega_ {2} = d x _ {2}, \quad \omega_ {3} = - x _ {1} d x _ {2} + d x _ {3}
\]

satisfy the relations

\[
d \omega_ {k} = - \sum_ {i <   j} c _ {i j k} \omega_ {i} \wedge \omega_ {j} \quad (k = 1, 2, 3)
\]

and are linearly independent at each point of \(\mathbf{R}^3\) .

\begin{enumerate}
\def\labelenumi{(\alph{enumi})}
\setcounter{enumi}{1}
\tightlist
\item
  Deduce that there exists, on an open neighbourhood of 0 in \(\mathbf{R}^3\) , a Lie group germ structure \(\mathbf{G}\) such that \(\mathrm{L}(\mathrm{G}) = \mathfrak{g}\) and
\end{enumerate}

\[
(a _ {1}, a _ {2}, a _ {3}) (x _ {1}, x _ {2}, x _ {3}) = (x _ {1} + a _ {1}, x _ {2} + a _ {2}, x _ {3} + a _ {3} + a _ {1} x _ {2})\tag{1}
\]

for \((a_{1}, a_{2}, a_{3})\) , \((x_{1}, x_{2}, x_{3})\) sufficiently close to 0. Show that in fact formulae (1) define a nilpotent Lie group structure on \(\mathbf{R}^{3}\) .

¶ 6. Let G be a finite-dimensional Lie group, g its Lie algebra, g* the dual vector space of g, σ the adjoint representation of G on g and ρ the representation of G on g* defined by ρ(g) = tσ(g)-1 for all g ∈ G.

\begin{enumerate}
\def\labelenumi{(\alph{enumi})}
\item
  \(\mathrm{L}(\rho)(x) = -^{t}(\mathrm{ad} x)\) for all \(x \in g\) .
\item
  Let \(f \in \mathfrak{g}^*\) . Let \(G_f\) denote the stabilizer of \(f\) in \(G\) ; it is a Lie subgroup of \(G\) and \(\mathfrak{g}_f = L(G_f)\) is the set of \(x \in G\) such that \({}^t (\operatorname{ad}x)f = 0\) .
\item
  For \(x, y\) in \(\mathfrak{g}\) , we write \(\mathrm{B}_f(x, y) = f([x, y])\) . Show that \(\mathrm{B}_f\) is an alternating bilinear form on \(\mathfrak{g}\) , that the mapping \(s_f\) of \(\mathfrak{g}\) into \(\mathfrak{g}^*\) left associated with \(\mathrm{B}_f\) is the mapping \(x \mapsto {}^t (\mathrm{ad}x)\) \(f\) and that the orthogonal of \(\mathfrak{g}\) with respect to \(\mathrm{B}_f\) is \(\mathfrak{g}_f\) . Let \(\beta_f\) denote the non-degenerate alternating bilinear form on \(\mathfrak{g} / \mathfrak{g}_f\) (which is therefore of even dimension) derived from \(\mathrm{B}_f\) when passing to the quotient. We denote by \(\beta_f'\) the inverse form of \(\mathrm{B}_f\) on \((\mathfrak{g} / \mathfrak{g}_f)*\) .
\item
  Let \(\Omega\) denote an orbit of \(G\) on \(\mathfrak{g}^*\) ; we give it the manifold structure derived from that on \(G / G_f\) by transport of structure, where \(f\) is an arbitrary point of \(\Omega\) . Then \(G\) operates analytically on \(\Omega\) on the left. Let \(D\) be the associated law of infinitesimal operation. For all \(a \in \mathfrak{g}\) and all \(f \in \Omega\) ,
\end{enumerate}

\[
\mathrm{D} _ {a} (f) = - (^ {t} \mathrm{ad} a) f.
\]

The tangent subspace \(\mathrm{T}_{f}(\Omega)\) is the image of \(s_{f}\) , that is the orthogonal of \(g_{f}\) in \(g^{*}\) , and is canonically identified with \((\mathfrak{g}/\mathfrak{g}_{f})^{*}\) , so that \(s_{f}\) defines a canonical isomorphism \(r_{f}\) of \(\mathrm{T}_{f}(\Omega)\) onto its dual. The field \(f \mapsto \beta_{f}'\) is an analytic differential form \(\omega\) of degree 2 on \(\Omega\) , which is invariant under G. For \(f \in \Omega\) , \(a \in g\) , \(b \in g\) , \(\omega(\mathrm{D}_{a}(f), \mathrm{D}_{b}(f)) = f([a, b])\) .

\begin{enumerate}
\def\labelenumi{(\alph{enumi})}
\setcounter{enumi}{4}
\item
  Show that \(d\omega = 0\) .
\item
  If \(\alpha\) is a differential form of degree 1 on \(\Omega\) , the isomorphisms \(r_f\) allow us to identify \(\alpha\) with a vector field. Let A be the set of analytic functions on \(\Omega\) with values in K. For \(\phi, \psi\) in A, we write \([\phi, \psi] = \omega(d\phi, d\psi) \in A\) . Show that A has thus a Lie algebra structure.
\item
  For all \(a \in \mathfrak{g}\) , let \(\psi_{a}\) be the function \(f \mapsto f(a)\) on \(\Omega\) . Show that \(a \mapsto \psi_{a}\) is a homomorphism of \(\mathfrak{g}\) into the Lie algebra A and that the form \(d\psi_{a}\) is identified, by means of the isomorphisms \(r_{f}\) , with the vector field \(\mathrm{D}_{a}\) .
\item
  Let U be an open subset of \(\mathfrak{g}^*\) and \(\phi\) an analytic function on U. For all \(f\in \mathrm{U}\) , the differential \(d_{f}\phi\) of \(\phi\) at \(f\) is identified with an element of \(\mathfrak{g}\) . Suppose that \(\mathrm{X}\phi = 0\) for every vector field X defined by the action of G on \(\mathfrak{g}^*\) . Show that then, for all \(f\in \mathrm{U}\) , \(d_{f}\phi\) belongs to the centre of \(\mathfrak{g}_f\) .
\item
  For all \(f \in \mathfrak{g}^*\) , we write \(r_f = \dim \mathfrak{g}_f\) . Let \(r = \inf_{f \in \mathfrak{g}^*} r_f\) . Show that the set \(V\) of \(f \in \mathfrak{g}^*\) such that \(r_f = r\) is open and dense in \(\mathfrak{g}^*\) . Show that, for all \(f \in V\) , \(\mathfrak{g}_f\) is commutative. (Construct \(r\) functions \(\phi\) satisfying the conditions of \((h)\) in a neighbourhood of \(f\) and such that their differentials at \(f\) are linearly independent.)
\end{enumerate}

¶ 7. (a) Let \((\lambda, x) \mapsto \lambda.x\) be a law of continuous left operation of R on T. Then one of the two following holds:

\begin{enumerate}
\def\labelenumi{(\arabic{enumi})}
\item
  either there exists a fixed point and the stabilizer of every point is either \(\mathbf{R}\) or \(\{0\}\) ;
\item
  or there exists a point whose stabilizer is an infinite discrete subgroup of \(\mathbf{R}\) and \(\mathbf{R}\) operates transitively on \(\mathbf{T}\) . (If \(\{\lambda \in \mathbf{R}|\lambda .x = x\} = \{0\}\) , the orbit of \(x\) is homeomorphic to \(\mathbf{R}\) and has frontier points \(\lim_{\lambda \to +\infty}\lambda .x\) which are fixed. If \(\{\lambda \in \mathbf{R}|\lambda .x = x\} = \mathbf{Z}a\) with \(a\neq 0\) , the orbit of \(x\) is homeomorphic to \(\mathbf{T}\) and \(\mathbf{R}\) is transitive.)
\end{enumerate}

\begin{enumerate}
\def\labelenumi{(\alph{enumi})}
\setcounter{enumi}{1}
\item
  If \(f\) is a homeomorphism of \(\mathbf{T}\) onto itself and \(k\) an integer \(\geqslant 1\) , a point \(x \in \mathbf{T}\) is called periodic of period \(k\) for \(f\) if \(f^k(x) = x\) and \(f^h(x) \neq x\) for \(1 \leqslant h < k\) . In the notation of (a), let \(f_{\lambda}\) be the homeomorphism \(x \mapsto \lambda .x\) . Then the set of periodic points of period \(k > 1\) for \(f_{\lambda}\) is empty or equal to \(\mathbf{T}\) . (If there exists a periodic point of period \(k > 1\) for \(f_{\lambda}\) , its stabilizer is distinct from \{0\} and \(\mathbf{R}\) , hence \(\mathbf{R}\) operates transitively on \(\mathbf{T}\) and all the points of \(\mathbf{T}\) are periodic of period \(k\) for \(f_{\lambda}.\) )
\item
  Let \(\Omega\) be a neighbourhood of \(e\) in \(\mathrm{Diff}^{\infty}(\mathbf{T})\) (Differentiable and Analytic Manifolds, R, 15.3.8, Note (1)). There exist an integer \(k > 1\) and \(f\in \Omega\) without fixed point such that the set of periodic points of period \(k\) for \(f\) is neither empty nor equal to \(\mathbf{T}\) . (Let \(p\colon \mathbf{R}\to \mathbf{T}\) be the canonical mapping. Let \(\rho\) be the mapping \(p(x) \mapsto p\left(x + \frac{2\pi}{k}\right)\) of \(\mathbf{T}\) into \(\mathbf{T}\) . Let \(\Omega'\) be a neighbourhood of \(e\) in \(\text{Diff}^\infty(\mathbf{T})\) such that \(\Omega'^2 \subset \Omega\) . For sufficiently large \(k, \rho \in \Omega'\) . On the other hand, let \(\sigma\) be a mapping of \(\mathbf{T}\) into \(\mathbf{T}\) such that \(\sigma(y) = y\) for
\end{enumerate}

\[
y \notin p \left(\right] 0, \frac {2 \pi}{k} [)
\]

and \(\sigma(p(x)) = p(x + \lambda(x))\) for \(x \in \left]0, \frac{2\pi}{k}\right\}\) , where

\[
0 <   \lambda (x) <   \frac {2 \pi}{k} - x;
\]

can be chosen such that \(\sigma \in \Omega'\) . Then \(f = \rho \circ \sigma\) has the required properties.) Such an \(f\) cannot, by (b), be in the image of a continuous morphism of \(\mathbf{R}\) into \(\operatorname{Diff}^{\infty}(\mathbf{T})\) .

\begin{enumerate}
\def\labelenumi{(\alph{enumi})}
\setcounter{enumi}{3}
\item
  Let Y be a compact differentiable manifold of dimension \(\geqslant1\) . Every neighbourhood of e in \(\operatorname{Diff}^{\infty}(Y)\) contains an element h with the following property: h belongs to the image of no continuous morphism of R into \(\operatorname{Diff}^{0}(Y)\) . (The manifold Y contains an open submanifold of the form \(T \times D\) , where D is the open Eulcidean ball of radius 1 and centre 0 in \(R^{n}\) ( \(n = \dim Y - 1\) ). Let \(g \in \operatorname{Diff}^{\infty}(D)\) be such that \(g(0) = 0\) , \(\|g(x)\| > \|x\|\) for \(0 < \|x\| < \frac{1}{2}\) and \(g(x) = x\) for \(\|x\| \geqslant \frac{1}{2}\) and g very close to e in \(\operatorname{Diff}^{\infty}(D)\) . Let f be as in (e), and, for 0 \textless{} t \textless{} 1, let \(f_{i}: T \to T\) be defined as follows: for \(x \in T\) , \(y \in p^{-1}(x)\) , \(z \in p^{-1}(f(x))\) and \(|z - y| < \pi\) , we write \(f_{i}(x) = p(ty + (1 - t)z)\) . Then there exists \(h \in \operatorname{Diff}^{\infty}(Y)\) such that \(h(x, y) = (f_{2\|x\|}(x), g(y))\) for \(x \in T\) , \(y \in D\) and \(h(u) = u\) for \(u \in Y - (T \times D)\) ; moreover, we can choose f and g such that h is arbitrarily close to e in \(\operatorname{Diff}^{\infty}(Y)\) . The points of \(T \times \{0\}\) are the \(y \in Y\) such that when n tends to +∞, \(h^{kn}(x)\) tends to a point not fixed under h (in fact it is periodic of period k). Therefore, every homeomorphism of Y onto Y which commutes with h leaves \(T \times \{0\}\) invariant. Then apply (e).
\item
  Let Y be a compact differentiable manifold of dimension \(\geqslant1\) . There do not exist a Lie group G and continuous morphisms \(\operatorname{Diff}^{\infty}(Y)\to G\) , \(G\to\operatorname{Diff}^{0}(Y)\) whose composition is the canonical injection of \(\operatorname{Diff}^{\infty}(Y)\) into \(\operatorname{Diff}^{0}(Y)\) . (Use (d).)
\end{enumerate}

\begin{enumerate}
\def\labelenumi{\arabic{enumi}.}
\setcounter{enumi}{7}
\tightlist
\item
  Let G be a finite-dimensional Lie group. Suppose that L(G) is simple. Let A be a normal subgroup of G. If A is not open, A is discrete and its centralizer in G is open. (Let a be the subalgebra tangent to A at e. Show that \(a = \{0\}\) . Let x be an element of A not belonging to the centre of G. Consider the mapping \(y \mapsto xyj^{-1}\) of G into A. Deduce from the relation \(a = \{0\}\) that the centralizer of x in G is open. Then, using an exponential mapping, show that the centralizer of x in G contains a neighbourhood of e independent of x.)
\end{enumerate}

¶ 9. Let G be a compact Lie group over K of dimension n.~Suppose that the Lie algebra L(G) is simple. For all \(g \in G\) , every integer \(m \geqslant 1\) and every neighbourhood V of e in G, let M(g, m, V) denote the set of elements of G of the form \(\prod_{i=1} x_i(y_i, g) x_i^{-1}\) with \(x_1, \ldots, x_m, y_1, \ldots, y_m\) in V.

\begin{enumerate}
\def\labelenumi{(\alph{enumi})}
\tightlist
\item
  Show that if \(g\) is an element of \(G\) whose centralizer is not open, if \(m\) is an integer \(\geqslant n\) and if \(V\) is a neighbourhood of \(e\) , then \(M(g, m, V)\) is a neighbourhood of \(e\) . (There exists \(a \in L(G)\) such that \((\operatorname{Ad} g)(a) \neq a\) . Let \(\phi\) be an exponential mapping of \(G\) . Let \(y\) be the image of 1 under the tangent mapping at 0 to \(\lambda \mapsto \langle \phi(xa), g \rangle\) (where \(\lambda \in K\) ). Then \(b = (\operatorname{Ad} g - 1)(a) \neq 0\) . There exist \(x_1, \ldots, x_m\) in \(V\) such that \(\{(\operatorname{Ad} x_1)(b), \ldots, (\operatorname{Ad} x_m)(b)\}\) contains a basis of \(L(G)\) . Consider the mapping
\end{enumerate}

\[
\left(x _ {1} ^ {\prime}, \dots , x _ {m} ^ {\prime}, y _ {1}, \dots , y _ {m}\right) \mapsto \prod_ {i = 1} ^ {n} x _ {i} ^ {\prime} (y _ {i}, g) x _ {i} ^ {\prime - 1}
\]

of \(\mathbf{G}^{2m}\) into G. Show that the tangent mapping at \((x_{1},\ldots ,x_{m},e,\ldots ,e)\) is surjective.)

\begin{enumerate}
\def\labelenumi{(\alph{enumi})}
\setcounter{enumi}{1}
\item
  Let \(\mathbf{U}\) be a neighbourhood of \(e\) in \(\mathbf{G}\) . Let \(m\) be an integer \(\geqslant 1\) . There exists an element \(g\) of \(\mathbf{G}\) whose centralizer is not open, such that \(\mathbf{M}(g, m, \mathbf{G}) \subset \mathbf{U}\) . (Argue by reductio ad absurdum using the compactness of \(\mathbf{G}\) .)
\item
  Let \(\mathbf{G}'\) be a compact Lie group over \(\mathbf{K}\) whose Lie algebra is simple. Let \(\phi: \mathbf{G} \to \mathbf{G}'\) be a bijective homomorphism of abstract groups. Then \(\phi\) is continuous. (Apply (b) to \(\mathbf{G}'\) and (a) to \(\mathbf{G}\) .)
\end{enumerate}

\begin{enumerate}
\def\labelenumi{\arabic{enumi}.}
\setcounter{enumi}{9}
\item
  Let \(G\) be a Lie group over \(K\) . Show that there exists a neighbourhood \(V\) of \(e\) with the following property: for every sequence \((x_{n})\) of elements of \(V\) , if we define inductively \(y_{1} = x_{1}, y_{n} = (x_{n}, y_{n-1})\) , the sequence \((y_{n})\) tends to \(e\) .
\item
  \begin{enumerate}
  \def\labelenumii{(\alph{enumii})}
  \tightlist
  \item
    For \(x, y\) in \(\mathbf{S}_{2n-1} \subset \mathbf{C}^n\) , we define \(\alpha(x, y) \in [0, \pi]\) by
  \end{enumerate}
\end{enumerate}

\[
\cos \alpha (x, y) = \mathscr {R} ((x | y)).
\]

For \(s, t\) in \(\mathbf{U}(n)\) , we write \(d(s, t) = \sup_{x \in \mathbf{S}_{2n-1}} \alpha(sx, tx)\) . Show that \(d\) is a left and right invariant distance on \(\mathbf{U}(n)\) .

\begin{enumerate}
\def\labelenumi{(\alph{enumi})}
\setcounter{enumi}{1}
\item
  Let \(s \in \mathbf{U}(n)\) . Let \(\theta_1, \ldots, \theta_m\) be the numbers in \(\mathbf{J} - \pi, \pi \mathbf{j}\) such that the \(e^{i\theta_j}\) are the distinct eigenvalues of \(s\) . Let \(\theta(s) = \sup_{1 \leqslant j \leqslant m} |\theta_j|\) . Then \(d(e, s) = \theta(s)\) . (Let \(V_j\) be the eigenspace of \(s\) corresponding to \(e^{i\theta_j}\) . For \(x \in S_{2n-1}\) , find a lower bound to \(\mathcal{R}((x|sx))\) by decomposing with respect to the \(V_j\) .)
\item
  Let \(s, t\) be in \(\mathbf{U}(n)\) be such that \(\theta(t) < \frac{\pi}{2}\) and \(s\) commutes with \((s, t)\) . Then \(s\) and \(t\) commute. (In the notation of (b), let \(\mathrm{V}_j'\) be the orthogonal supplement of \(\mathrm{V}_j\) and \(\mathrm{W}_j = t(\mathrm{V}_j)\) ; show that \(\mathrm{W}_j = (\mathrm{W}_j \cap \mathrm{V}_j) + (\mathrm{W}_j \cap \mathrm{V}_j')\) and then that \(\mathrm{W}_j \cap \mathrm{V}_j' = \{0\}\) .)
\item
  Show that there exists a compact neighbourhood \(\mathbf{V}\) of \(e\) in \(\mathbf{U}(n)\) which has the property of Exercise 10, which is stable under inner automorphisms of \(\mathbf{U}(n)\) and which is such that if \(x, y\) are non-permutable elements of \(\mathbf{V}\) then \(x\) and \((x, y)\) are not permutable. (Use (c).)
\item
  Let \(\beta\) be the normalized Haar measure of \(\mathbf{U}(n)\) . Let \(\mathrm{V}_1\) be a symmetric compact neighbourhood of \(e\) in \(\mathbf{U}(n)\) such that \(\mathrm{V}_1^2 \in \mathrm{V}\) . Let \(p\) be an integer such that \(\beta (\mathrm{V}_1) > 1 / p\) . Show that, for every finite subgroup F of \(\mathbf{GL}(n,\mathbf{C})\) , there exists a commutative normal subgroup A of F such that \(\operatorname{Card}(\mathrm{F} / \mathrm{A}) \leqslant p\) (Jordan's theorem). (Reduce it to the case where \(\mathrm{F} \subset \mathbf{U}(n)\) . Take A to be the subgroup of F generated by \(\mathrm{F} \cap \mathrm{V}\) and use (d).)
\end{enumerate}

\begin{enumerate}
\def\labelenumi{\arabic{enumi}.}
\setcounter{enumi}{11}
\tightlist
\item
  Let G be a finite-dimensional connected real or complex Lie group and \(\alpha\) a left invariant differential form of degree p on G. For \(\alpha\) to be right invariant too, it is necessary and sufficient that \(d\alpha = 0\) . (Observe that the condition for right invariance can be written
\end{enumerate}

\[
\wedge^ {p} (^ {t} \mathrm{Ad} (s)). \alpha (e) = \alpha (e)
\]

for all \(s \in G\) and hence it is equivalent to the condition

\[
\sum_ {j = 1} ^ {f} \left\langle \alpha (e), u _ {1} \wedge \dots \wedge u _ {j - 1} \wedge [ u, u _ {j} ] \wedge u _ {j + 1} \wedge \dots \wedge u _ {p} \right\rangle = 0
\]

for all \(u, u_1, \ldots, u_p\) in \(\mathrm{L}(\mathrm{G})\) . In particular, the left and right invariant differential forms of degree 1 on \(\mathbf{G}\) form a vector space of dimension

\[
\dim \mathrm{L} (\mathrm{G}) - \dim [ \mathrm{L} (\mathrm{G}), \mathrm{L} (\mathrm{G}) ].
\]

\begin{enumerate}
\def\labelenumi{\arabic{enumi}.}
\setcounter{enumi}{12}
\tightlist
\item
  Let \(\mathbf{R}^n\) be given the usual scalar product \(((\xi_i), (\eta_i)) \mapsto \sum_{i=1} \xi_i \eta_i\) . Let \(\mathbf{I}(n)\) be the group of isometric affine transformations of \(\mathbf{R}^n\) . It is canonically isomorphic to the Lie subgroup of \(\mathbf{GL}(n + 1, \mathbf{R})\) consisting of the matrices
\end{enumerate}

\[
S = \left( \begin{array}{c c} U & x \\ 0 & 1 \end{array} \right)
\]

where \(U \in \mathbf{O}(n)\) and \(x\) is an arbitrary element of \(\mathbf{R}^n\) (matrix of type \((n,1)\) ). \(\mathbf{I}(n)\) can be identified with the semi-direct product of \(\mathbf{O}(n)\) by \(\mathbf{R}^n\) defined by the canonical injection of \(\mathbf{O}(n)\) into \(\mathbf{GL}(n,\mathbf{R})\) (the matrix \(S\) is then denoted by \((U,x)\) ). Let \(p\colon \mathbf{I}(n)\to \mathbf{O}(n)\) be the canonical surjection, with kernel \(\mathbf{R}^n\) .

\begin{enumerate}
\def\labelenumi{(\alph{enumi})}
\tightlist
\item
  Let \(\Gamma\) be a discrete subgroup of \(\mathrm{I}(n)\) . Consider the closure \(\tilde{p}(\Gamma)\) of \(p(\Gamma)\) in \(\mathbf{O}(n)\) . Show that its identity component is commutative. (Let \(V\) be a neighbourhood of \(I\) in \(\mathbf{O}(n)\) with the properties of Exercise 11 ( \(d\) ) and such that further \(\|U - I\| \leqslant \frac{1}{4}\) for all \(U \in V\) (the norm on \(\operatorname{End}(\mathbf{R}^n)\) being derived from the Euclidean norm on \(\mathbf{R}^n\) ). Argue by reductio ad absurdum, assuming that there exist \(S_1 = (U_1, x_1)\) , \(S_2 = (U_2, x_2)\) such that \(U_1\) and \(U_2\) belong to \(V\) and do not commute. Show that \((S_1, S_2) = ((U_1, U_2), y)\) with
\end{enumerate}

\[
\| y \| \leqslant \frac {1}{4} (\| x _ {1} \| + \| x _ {2} \|).
\]

Then define the \(S_{k}\) inductively by \(S_{k}=(S_{1},S_{k-1})\) for \(k\geqslant3\) . Using the choice of V, show that this sequence has an infinity of distinct terms and is bounded in \(\mathbf{I}(n)\) , which is absurd.)

\begin{enumerate}
\def\labelenumi{(\alph{enumi})}
\setcounter{enumi}{1}
\item
  \(\Gamma\) is called crystallographic if \(\mathrm{I}(n)/\Gamma\) is compact. Show that if this is so, for all \(x \in \mathbf{R}^n\) , the affine subspace \(\mathbf{L}\) of \(\mathbf{R}^n\) generated by \(\Gamma x\) is equal to \(\mathbf{R}^n\) . (Otherwise, for all \(y \in \mathbf{R}^n\) , all the points of \(\Gamma y\) are the same distance from \(\mathbf{L}\) and, as this distance can be arbitrarily large \(\mathrm{I}(n)/\Gamma\) is not compact.)
\item
  If \(\Gamma\) is commutative and crystallographic, then \(\Gamma \subset \mathbf{R}^n\) . (If \(S = (U, x) \in \Gamma\) is such that the vector subspace \(\mathrm{V} \subset \mathbf{R}^n\) consisting of the points invariant under \(U\) is not equal to \(\mathbf{R}^n\) and \(\mathrm{V}^\perp\) denotes the orthogonal subspace of \(\mathrm{V}\) , then, for all \(S' = (U', x') \in \Gamma\) , \(U'\) leaves \(\mathrm{V}\) and \(\mathrm{V}^\perp\) stable. On the other hand, since \(U|\mathrm{V}^\perp\) has no non-zero invariant point, show that this allows us to assume, by changing the origin, that \(x \in \mathrm{V}\) . Using (b), there exists \(S' = (U', x') \in \Gamma\) such that the orthogonal projection \(y'\) of \(x'\) onto \(\mathrm{V}^\perp\) is \(\neq 0\) . By evaluating \(Sx'\) by the formula \(S = S'SS'^{-1}\) , deduce that \(U.y' = y'\) contrary to the definition of \(\mathrm{V}\) .
\item
  If \(\Gamma\) is crystallographic, \(\Gamma \cap \mathbf{R}^n\) is a free commutative group of rank \(n\) and \(\Gamma / (\Gamma \cap \mathbf{R}^m)\) is a finite group (Bieberbach's Theorem). (Let \(W\) be the vector subspace of \(\mathbf{R}^n\) generated by \(\Gamma \cap \mathbf{R}^n\) . The compact group \(p(\Gamma)\) has only a finite number of connected components. If \(W = \{0\}\) , \(\Gamma\) contains, by \((a)\) , a commutative normal subgroup of finite index \(\Gamma_1\) ; \(\Gamma_1\) is then crystallographic, which implies a contradiction with \((c)\) . Hence \(W \neq \{0\}\) . Show that \(p(\Gamma)\) leaves \(W\) stable and that \(p(\Gamma)|W\) is finite; otherwise there would exist a sequence \((S_m)\) in \(\Gamma\) such that if \(S_m = (U_m, x_m)\) , the \(U_m\) are distinct and tend to \(I\) ; then form the \((I, a_j)S_m(I, a_j)^{-1s_m^1}\) , where \((a_j)\) is a basis of \(\Gamma \cap \mathbf{R}^n\) , and obtain a contradiction with the hypothesis that \(\Gamma\) is discrete. Finally, to see that \(W = \mathbf{R}^n\) , show that otherwise the action of \(\Gamma\) on \(\mathbf{R}^n/W\) would be that of a crystallographic group containing no translation \(\neq 0\) .)
\end{enumerate}

